\section{Preliminaries}

\subsection{Fusion 2-categories}

Over an algebraically closed field $\mathds{k}$ of characteristic zero, the notions of finite semisimple 2-category and fusion 2-category were introduced in \cite{DR}. We succinctly review these definitions here, and refer the reader to the aforementioned reference and also \cite{D2} for details.

A ($\mathds{k}$-linear) 2-category $\mathfrak{C}$ is locally finite semisimple if its $\operatorname{Hom}$-1-categories are finite semisimple. An object $C$ of $\mathfrak{C}$ is called simple if $\operatorname{Id}_C$ is a simple object of the finite semisimple 1-category~$\operatorname{End}_{\mathfrak{C}}(C)$. Then, a \textit{finite semisimple} 2-category is a locally finite semisimple 2-category that is Cauchy complete in the sense of \cite{GJF} (see also \cite{D1}), has right and left adjoints for 1-morphisms, and has finitely many equivalence classes of simple objects. In a finite semisimple 2-category $\mathfrak{C}$, any two simple objects are called connected if there exists a non-zero 1-morphism between them. This defines an equivalence relation, and we write $\pi_0(\mathfrak{C})$ for the corresponding set of equivalences classes; This is the set of \textit{connected components} of $\mathfrak{C}$.

Now, the definition of a monoidal 2-category is well known. In particular, using the notations of \cite{SP}, a monoidal structure on the 2-category $\mathfrak{C}$ involves a 2-functor $\Box\colon \mathfrak{C}\times \mathfrak{C}\rightarrow\mathfrak{C}$, the monoidal product, and a distinguished object $I$, the monoidal unit, together with various other coherence data. Given any object $C$ of $\mathfrak{C}$, we can ask for the existence of its left dual $^{\sharp}C$ and its right dual~$C^{\sharp}$. These two notions were studied in \cite{Pstr}. Categorifying the concept of a fusion 1-category is therefore straightforward, and we obtain the following definition.

\begin{Definition}
A multifusion 2-category is a finite semisimple monoidal 2-category $\mathfrak{C}$ such that every object $C$ admits a right dual $C^{\sharp}$ and a left dual $^{\sharp}C$. A fusion 2-category is a multifusion 2-category whose monoidal unit $I$ is a simple object.
\end{Definition}

To every fusion 2-category $\mathfrak{C}$, there is an associated braided fusion 1-category $\Omega\mathfrak{C}:=\operatorname{End}_{\mathfrak{C}}(I)$, the 1-category of endomorphisms of the monoidal unit $I$ of $\mathfrak{C}$. Conversely, given any braided fusion 1-category $\mathcal{B}$, we can consider the fusion 2-category $\mathbf{Mod}(\mathcal{B})$ of finite semisimple left $\mathcal{B}$-module 1-categories with monoidal structure given by $\boxtimes_{\mathcal{B}}$, the relative Deligne tensor product. Below, we give another class of examples that will be relevant for our purposes. Many more may be found in \cite{DR} and \cite{DY23}, and others will be discussed subsequently.

\begin{Example}
Given $G$ a finite group, and $\pi$ a 4-cocycle for $G$ with coefficients in $\mathds{k}^{\times}$, we can then consider the fusion 2-category $\mathbf{2Vect}^{\pi}(G)$ of $G$-graded finite semisimple 1-categories, also known as 2-vector spaces, with pentagonator twisted by $\pi$. Such fusion 2-category were completely characterized in \cite{JFY} as those fusion 2-categories $\mathfrak{C}$ for which $\Omega\mathfrak{C}\simeq \mathbf{Vect}$. We call these bosonic strongly fusion 2-categories.
\end{Example}

\subsection{Module 2-categories and the relative 2-Deligne tensor product}

Let $\mathfrak{C}$ be a monoidal 2-category with monoidal product $\Box$, and monoidal unit $I$. A right $\mathfrak{C}$-module 2-category is a 2-category $\mathfrak{M}$ equipped with an action 2-functor $\Box\colon\mathfrak{M}\times\mathfrak{C}\rightarrow \mathfrak{M}$, and coherence data satisfying various axioms. Likewise, there is a notion of left $\mathfrak{C}$-module 2-category. Additionally, there are notions of $\mathfrak{C}$-module 2-functors, $\mathfrak{C}$-module 2-natural transformations, and~$\mathfrak{C}$-module modifications. There are also notions of bimodule 2-categories, and maps between them. We refer the reader to \cite[Section~2]{D4} for the precise definitions.

Let us now assume that $\mathfrak{C}$ is a fusion 2-category. If $\mathfrak{M}$ is a left $\mathfrak{C}$-module 2-category, then we can consider $\mathbf{\operatorname{End}}_{\mathfrak{C}}(\mathfrak{M})$, the monoidal 2-category of left $\mathfrak{C}$-module 2-endofunctors on $\mathfrak{M}$. The next result combines \cite[Theorem~5.3.2]{D8} with \cite[Corollary~5.1.3]{D9}.

\begin{Theorem}\label{thm:Moritadual}
The monoidal $2$-category $\mathbf{\operatorname{End}}_{\mathfrak{C}}(\mathfrak{M})$ is a multifusion $2$-category.
\end{Theorem}

Moreover, the multifusion 2-category $\mathbf{\operatorname{End}}_{\mathfrak{C}}(\mathfrak{M})$ is fusion if and only if $\mathfrak{M}$ is indecomposable as a finite semisimple left $\mathfrak{C}$-module 2-category.

For our purposes, we need to consider another operation on finite semisimple module 2-categories. Let $\mathfrak{M}$ be a finite semisimple right $\mathfrak{C}$-module 2-category and $\mathfrak{N}$ be a finite semisimple left $\mathfrak{C}$-module 2-category. There are notions of $\mathfrak{C}$-balanced 2-functors, $\mathfrak{C}$-balanced 2-natural transformations, and $\mathfrak{C}$-balanced modifications out of $\mathfrak{M}\times\mathfrak{N}$. We refer the reader to \cite[Section~2.2]{D10} for the details. The next result is a combination of \cite[Theorem~2.2.4]{D10} with \cite[Corollary~5.1.3]{D9}.

\begin{Theorem}
There is a $3$-universal $\mathfrak{C}$-balanced $2$-functor $\boxtimes_{\mathfrak{C}}\colon \mathfrak{M}\times\mathfrak{N}\rightarrow \mathfrak{M}\boxtimes_{\mathfrak{C}}\mathfrak{N}$ to a finite semisimple $2$-category.
\end{Theorem}

In fact, the proof of \cite[Theorem~2.2.4]{D10} gives an explicit description of $\mathfrak{M}\boxtimes_{\mathfrak{C}}\mathfrak{N}$. More precisely, if $A$ is a (necessarily separable) algebra in $\mathfrak{C}$ such that $\mathfrak{M}$ is equivalent to $\mathbf{LMod}_{\mathfrak{C}}(A)$, the right~$\mathfrak{C}$-module 2-category of left $A$-modules in $\mathfrak{C}$, and $B$ is a (necessarily separable) algebra in $\mathfrak{C}$ such that $\mathfrak{N}$ is equivalent to $\mathbf{Mod}_{\mathfrak{C}}(B)$, the left $\mathfrak{C}$-module 2-category of right $B$-modules in~$\mathfrak{C}$, then~$\mathfrak{M}\boxtimes_{\mathfrak{C}}\mathfrak{N}\simeq \mathbf{Bimod}_{\mathfrak{C}}(A,B)$, the finite semisimple 2-category of $A$-$B$-bimodules in $\mathfrak{C}$.

Finally, let us recall that, as was explained in \cite[Section~3.2]{D10}, the existence of the relative 2-Deligne tensor product allows us to consider the symmetric monoidal Morita 4-category $\mathbf{F2C}$ of (multi)fusion 2-categories, finite semisimple bimodule 2-categories, and bimodule morphisms.

\section{Extension theory}

For fusion 1-categories over an algebraically closed field of characteristic zero, extension theory was developed in \cite{ENO2}. The key concept is that of an invertible bimodule 1-category, or, more precisely, the space formed by such objects together with their invertible morphisms. Proceeding in a similar fashion, we will discuss extension theory for fusion 2-categories, i.e., we will study group graded fusion 2-categories in the sense of the definition below.

\begin{Definition}
Let $\mathfrak{C}$ be a fusion 2-category and $G$ a finite group. A $G$-grading on $\mathfrak{C}$ is a~decomposition $\boxplus_{g\in G}\mathfrak{C}_g$ of $\mathfrak{C}$ into a direct sum of finite semisimple 2-categories such that for every $C$ in $\mathfrak{C}_g$ and $D$ in $\mathfrak{C}_h$, $C\Box D$ lies $\mathfrak{C}_{gh}$. A $G$-grading on $\mathfrak{C}$ is faithful if $\mathfrak{C}_g$ is non-zero for every $g$ in $G$.
\end{Definition}

\subsection{Invertible bimodule 2-categories}

Let $\mathds{k}$ be an algebraically closed field of characteristic zero. We begin by recalling a definition from \cite{DY23}. Let $\mathfrak{C}$ and $\mathfrak{D}$ be two fusion 2-categories. We write $\mathfrak{D}^{\rm mop}$ for the fusion 2-category obtained by endowing the finite semisimple 2-category $\mathfrak{D}$ with the opposite of the monoidal product of $\mathfrak{D}$.

\begin{Definition}
A finite semisimple $\mathfrak{C}$-$\mathfrak{D}$-bimodule 2-category $\mathfrak{M}$ is invertible if the canonical monoidal 2-functor $\mathfrak{D}^{\rm mop}\rightarrow \mathbf{\operatorname{End}}_{\mathfrak{C}}(\mathfrak{M})$ is an equivalence.
\end{Definition}

At the decategorified level, i.e., for fusion 1-categories, invertibility of a finite semisimple bimodule 1-category admits various equivalent characterizations as is explained in \cite[Proposition~4.2]{ENO2}. A categorified version of this result was given in \cite{D10}, which we partially recall below.

\begin{Proposition}\label{prop:invertibilitycharacterization}
Let $\mathfrak{M}$ be a finite semisimple $\mathfrak{C}$-$\mathfrak{D}$-bimodule $2$-category. The following are equivalent:
\begin{enumerate}\itemsep=0pt
 \item[$(1)$] The $\mathfrak{C}$-$\mathfrak{D}$-bimodule $2$-category $\mathfrak{M}$ is invertible.
 \item[$(2)$] The $\mathfrak{C}$-$\mathfrak{D}$-bimodule $2$-category $\mathfrak{M}$ defines an invertible $1$-morphism from $\mathfrak{C}$ to $\mathfrak{D}$ in $\mathbf{F2C}$.
\end{enumerate}
\end{Proposition}

The next proposition gives preliminary insight into the relation between group graded fusion 2-categories and invertible bimodule 2-categories. We note that the first part of the result below has already appeared as \cite[Proposition~3.1.7]{DY23}. We also refer the reader to \cite[Theorem~6.1]{ENO2} for the decategorified version of this proposition.

\begin{Proposition}\label{prop:gradedMoritaequivalence}
Let $\mathfrak{C}$ be a faithfully $G$-graded fusion $2$-category $\mathfrak{C}$. For any $g\in G$, the finite semisimple $\mathfrak{C}_e$-$\mathfrak{C}_e$-bimodule $2$-category $\mathfrak{C}_g$ is invertible. Furthermore, for any $g,h\in G$, the $2$-functor $\Box\colon\mathfrak{C}\times\mathfrak{C}\rightarrow \mathfrak{C}$ induces an equivalence \smash{$\mathfrak{C}_g\boxtimes_{\mathfrak{C}_e}\mathfrak{C}_h\xrightarrow{\simeq}\mathfrak{C}_{gh}$} of $\mathfrak{C}_e$-$\mathfrak{C}_e$-bimodule $2$-categories.
\end{Proposition}
\begin{proof}
The first part follows from the second, so we only prove the latter. For the second part, note that it is enough to show that the canonical 2-functor is an equivalence. To this end, pick any non-zero objects $X$ in $\mathfrak{C}_{g^{-1}}$, and $Y$ in $\mathfrak{C}_{h}$. It follows from the proof of \cite[Theorem~5.4.3]{D8} that~$Y\Box {^{\sharp}Y}$ is a separable algebra in $\mathfrak{C}$, and that the left $\mathfrak{C}$-module 2-functor $\mathfrak{C}\rightarrow \mathbf{Mod}_{\mathfrak{C}}\big(Y\Box {^{\sharp}Y}\big)$ given by $C\mapsto C\Box {^{\sharp}Y}$ is an equivalence. Likewise, $X\Box {^{\sharp}X}$ is a separable algebra in $\mathfrak{C}$, and the right~$\mathfrak{C}$-module 2-functor $\mathfrak{C}\rightarrow \mathbf{LMod}_{\mathfrak{C}}\big(X\Box {^{\sharp}X}\big)$ given by $C\mapsto X\Box C$ is an equivalence. In particular, we also find that the 2-functor $\mathfrak{C}\rightarrow \mathbf{Bimod}_{\mathfrak{C}}\big(X\Box {^{\sharp}X},Y\Box {^{\sharp}Y}\big)$ given by $C\mapsto X\Box C\Box {^{\sharp}Y}$ is an equivalence. Under these identifications, as was recalled above, it follows from the proof of~\cite[Theorem~2.2.4]{D10} that the canonical 2-functor $\mathfrak{C}\times\mathfrak{C}\rightarrow\mathfrak{C}\boxtimes_{\mathfrak{C}}\mathfrak{C}$ is identified with $\Box\colon\mathfrak{C}\times\mathfrak{C}\rightarrow \mathfrak{C}$, which is given by
\begin{gather*}
 \mathbf{LMod}_{\mathfrak{C}}\big(X\Box {^{\sharp}X}\big)\times\mathbf{Mod}_{\mathfrak{C}}\big(Y\Box {^{\sharp}Y}\big) \rightarrow \mathbf{Bimod}_{\mathfrak{C}}\big(X\Box {^{\sharp}X},Y\Box {^{\sharp}Y}\big),\qquad
(M,N) \mapsto M\Box N.
\end{gather*}
 But, by considering the restriction $\mathfrak{C}_e\hookrightarrow\mathfrak{C}$, the above equivalences restrict to equivalences $\mathfrak{C}_h\rightarrow \mathbf{Mod}_{\mathfrak{C}_e}\big(Y\Box {^{\sharp}Y}\big)$, $\mathfrak{C}_g\rightarrow \mathbf{LMod}_{\mathfrak{C}_e}\big(X\Box {^{\sharp}X}\big)$, and $\mathfrak{C}_{gh}\rightarrow \mathbf{Bimod}_{\mathfrak{C}_e}\big(X\Box {^{\sharp}X},Y\Box {^{\sharp}Y}\big)$. In particular, the canonical 2-functor $\mathfrak{C}_g\times\mathfrak{C}_h\rightarrow\mathfrak{C}_g\boxtimes_{\mathfrak{C}_e}\mathfrak{C}_h$ is identified with $\Box:\mathfrak{C}_g\times\mathfrak{C}_h\rightarrow\mathfrak{C}_{gh}$ as claimed.
\end{proof}

A similar argument as the one used in the proof of the above proposition yields the following result.

\begin{Corollary}\label{cor:invertibledual}
The canonical $\mathfrak{C}_e$-$\mathfrak{C}_e$-bimodule $2$-functor $\mathfrak{C}_{g^{-1}}\rightarrow \mathbf{Fun}_{\mathfrak{C}_e}(\mathfrak{C}_g,\mathfrak{C}_e)$ given by $D\mapsto \{C\mapsto D\Box C\}$ is an equivalence.
\end{Corollary}

\subsection{Brauer--Picard spaces and extensions}

Group-graded extensions of fusion 1-categories are parameterised by maps into the space of invertible bimodule 1-categories and their (invertible) higher morphisms \cite{ENO2}. Thanks to the existence of the 4-category $\mathbf{F2C}$ obtained in \cite{D10}, the corresponding spaces for our 2-categorical purposes are easy to define.

\begin{Definition}
Let $\mathfrak{C}$ be a fusion 2-category. The Brauer--Picard space of $\mathfrak{C}$ consists of the invertible objects and the invertible morphisms in the monoidal 3-category of finite semisimple $\mathfrak{C}$-$\mathfrak{C}$-bimodule 2-categories, that is, \begin{gather*}\mathscr{B}{\rm r}\mathscr{P}{\rm ic}(\mathfrak{C}):=\big(\mathbf{\operatorname{End}}_{\mathbf{F2C}}(\mathfrak{C})\big)^{\times}.\end{gather*}
\end{Definition}

\begin{Remark}
We will write $Br\operatorname{Pic}(\mathfrak{C}):= \pi_0(\mathscr{B}{\rm r}\mathscr{P}{\rm ic}(\mathfrak{C}))$, the Brauer--Picard group of $\mathfrak{C}$. It follows from \cite[Lemma~2.2.1]{D9} that $\Omega \mathbf{\operatorname{End}}_{\mathbf{F2C}}(\mathfrak{C})\simeq \mathscr{Z}(\mathfrak{C})$, the Drinfeld center of $\mathfrak{C}$, as defined in~\cite{Cr}. Thus, the homotopy groups of the space $\mathscr{B}{\rm r}\mathscr{P}{\rm ic}(\mathfrak{C})$ are given as follows:
 $$\renewcommand{\arraystretch}{1.2} \begin{tabular}{|c|c|c|c|c|}
\hline
$\pi_0$ & $\pi_1$ & $\pi_2$ & $\pi_3$\\
\hline
${\rm Br}\operatorname{Pic}(\mathfrak{C})$ & $\operatorname{Inv}(\mathscr{Z}(\mathfrak{C}))$ & $\operatorname{Inv}(\Omega\mathscr{Z}(\mathfrak{C}))$ & $\mathds{k}^{\times}$\\
\hline
\end{tabular}$$
\end{Remark}
